\paragraph{Source-qualified energies in GeV.}
A positive calibration $u$ sends a source Hamiltonian to
$H_{\rm GeV}=u(H-E_{\rm ref}I)$ and its clock to $\tau=t/u$.
The two propagators differ by the scalar phase $e^{itE_{\rm ref}}$;
therefore every complete apparatus probability, including preparation
failure, is invariant. The charged band minima are
$0.00051099895069$, $0.10565941449590446$ and $1.776984970813427$ GeV.
The electron supplies the scale, and the analyzer phase $2/9$ is a
retrospective architecture choice. A separately calibrated Dirac neutral
tensor gives $(1.0000000,8.7120606,50.1398045)\times10^{-12}$ GeV,
conditional on its supplied oscillation parameters, normal ordering and
lightest mass $0.001$ eV. Its active and sterile chiral coordinates share
these masses. The full quadratic comparison also supplies quark running
parameters and tree boson coefficients; these are kept distinct from
interacting physical rest masses. Complete baryon spectral windows retain
their preparation failure and fitted input ancestry.
The spin-gradient member fits the deuteron-assigned ground energy
near $1.87561294497$ GeV with its fifth empirical input, retaining all
unassigned bound channels. The positive-frequency weak member gives
$\Gamma_\mu=3.00876260\times10^{-19}$ GeV and a tau leptonic
width ratio $0.972563836$, compared with $0.9730\pm0.0022$.
The ratio is an unused comparison; the weak normalization has
muon-lifetime ancestry. The executable catalogue is
\texttt{model/openwave\_mass\_catalogue}.
